% === Revised Preface + Section 1 ===
%
% NEW PREAMBLE ADDITIONS NEEDED (lift these into the main preamble):
%
% \usepackage{tikz}
% \usetikzlibrary{positioning,arrows.meta,calc}
% \usepackage{booktabs}
% \usepackage{listings}
%
% % --- activity tcolorbox (green accent, distinct from orange "intuition") ---
% \newtcolorbox{activity}[1][]{
%   breakable,
%   enhanced,
%   colback=green!5,
%   colframe=green!45!black,
%   coltitle=white,
%   fonttitle=\bfseries,
%   title={Activity},
%   boxrule=0.6pt,
%   left=6pt,right=6pt,top=4pt,bottom=4pt,
%   #1
% }
%
% % --- "how to read" micro-box (grey, minimal) ---
% \newtcolorbox{howtoread}{
%   enhanced,
%   colback=gray!8,
%   colframe=gray!50,
%   boxrule=0.4pt,
%   left=6pt,right=6pt,top=3pt,bottom=3pt,
%   sharp corners,
%   fontupper=\small,
% }
%
% ---------------------------------------------------------------------------

% =========================================================================
\section*{Preface: how to read these notes}
% =========================================================================

\begin{howtoread}
\textbf{Two readers, one document.}
Orange \textbf{Intuition} boxes are the \emph{manager track}: plain
English, small numbers, one takeaway each. The running body is the
\emph{student track}: definitions, lemmas, proofs. Green
\textbf{Activity} boxes are for either reader --- they finish in under
ten minutes and produce a number you can argue about.
If you only read the orange boxes and the section headings, you will
still leave with the thesis.
\end{howtoread}

\paragraph{The thesis in one sentence.}
\emph{Predicting the next symbol well, compressing a file well, and
acting intelligently are --- under precise definitions given later ---
the same operation measured in different units.}
That is either the most important business fact of the decade, or a
confusion between three different things that happen to share a name.
By the end of these notes you will know which, and you will know what
observation would prove you wrong.

\tableofcontents
\newpage

% =========================================================================
\section{Why this equivalence matters}
\label{sec:why}
% =========================================================================

\subsection*{A story to hold in your head}

Three friends meet at a bar. The first is a stenographer: her job is
to guess, letter by letter, what the speaker will say next, and she
is eerily good at it. The second runs a file-compression startup: his
software turns a 1~GB text dump into a 150~MB archive without losing
a byte. The third is a chess player who has never read a book on
chess but wins anyway, because she has an uncanny sense of which move
a strong opponent would play.

Over the course of these notes we will prove that the three friends
are, in a precise mathematical sense, the same person. Shannon will
show the stenographer and the compressor are literally doing the same
job in different units. Kolmogorov will show that ``the shortest
description of a thing'' does not depend on which language you
describe it in. Solomonoff and Hutter (the latter an ETH alum, the
former a little-known American polymath --- we will introduce both
properly later) will show that a gambler who always bets on the
shortest explanation is, in a sense we will make exact, the smartest
possible learner. That is the whole plot.

\begin{intuition}
\textbf{One sentence.} Being good at predicting what comes next,
being good at compressing a file, and being good at thinking are ---
up to a translation we will spell out --- the same skill.

\smallskip
\textbf{Why a decision-maker should care.} Your vendors are selling
products whose entire training objective is one number: the average
surprise of the model when it sees the next word of the internet.
If the equivalence holds, minimising that one number is a proxy for
general intelligence. If it does not hold, you are buying a very
expensive autocomplete. You cannot tell which without seeing the
equivalence and its cracks.

\smallskip
\textbf{The falsifiable version.} If the equivalence is real, then
across language models of similar size, \emph{lower bits-per-
character on held-out English text should track higher scores on
reasoning benchmarks} (MMLU, BIG-Bench, GSM8K). If two models have
the same held-out bits-per-character but wildly different benchmark
scores, the equivalence is leaky and we should say where. This is
not a rhetorical flourish --- it is the empirical hook the rest of
the document hangs on.
\end{intuition}

\subsection*{A worked number to anchor the rest}

Before the definitions, here is the number to remember.
Uncompressed English Wikipedia is about 19~GB. The classical utility
\texttt{gzip} shrinks it to roughly 5~GB (a compression ratio of
about 3.8). A modestly tuned 7-billion-parameter language model, used
as a compressor through arithmetic coding, shrinks the same file to
under 3~GB --- a ratio above 6. The Hutter Prize, a long-running
competition founded by Marcus Hutter, currently stands near 1.1~GB
on a 1~GB benchmark slice, i.e.\ a ratio close to 9. The prize's
explicit premise is: \emph{the team that compresses English hardest
has, by that very fact, built the best model of English}.
\citep{hutter2006prize}

\begin{example}[Bits, dollars, and benchmark points]
\label{ex:three-currencies}
Suppose a firm stores 100~TB of customer-support transcripts at
\$0.023 per GB-month on commodity cloud storage. Raw cost:
about \$28{,}000/year. A foundation model that compresses the logs
2$\times$ (a realistic figure for natural-language logs) saves
roughly \$14{,}000/year. That is the \textbf{dollars column}.
The \textbf{bits column} is the bits-per-character of the same
model on a held-out slice of the same logs. The \textbf{capability
column} is what the same model scores on MMLU. The thesis of these
notes is that the three columns are, up to noise and the caveats of
\cref{sec:caveats}, three views of one underlying quantity.
\end{example}

\begin{activity}[title={Activity A1 --- Guess the next letter}]
\textbf{Materials.} A short sentence (10--20 words), one letter per
index card, kept face down.
\textbf{Protocol.} In pairs. Flip card 1. Before flipping card 2,
the guesser names the letter they think is next. Tally hits over 60
cards.
\textbf{Debrief.} Hit rate $h$ gives a back-of-envelope upper bound
on your per-letter surprise: roughly $-\log_2 h$ bits. A frontier LM
achieves about 0.9~bits per character on English. What did you
achieve? The gap between you and the LM, in bits, is the gap that
Shannon will soon teach us to price.
\end{activity}

\begin{activity}[title={Activity A2 --- Compress your own paragraph}]
\textbf{Snippet (Python, $\sim$5 lines).}
\begin{quote}\small\ttfamily
import gzip, sys\\
s = open("my\_paragraph.txt","rb").read()\\
print("gzip bpc:", 8*len(gzip.compress(s))/len(s))
\end{quote}
Run it on a paragraph you wrote yourself. Typical English prose
gives gzip $\approx 2.5$--$3.0$ bits per character. The same
paragraph scored by a small public LM (code omitted; see the course
repository) will give something closer to $1.0$. The ratio between
those two numbers is roughly the ratio of ``what a dumb model of
English knows'' to ``what a smart model of English knows.'' That
ratio, multiplied by your storage bill, is the dollar value of
intelligence measured in bits.
\end{activity}

\subsection*{The chain of results, informally}

The rest of the section previews, in English, the four technical
results that lock the equivalence together. Each will get a full
treatment later; here we only name them and say which gap they
close.

\begin{enumerate}[leftmargin=1.4em]
  \item \textbf{Shannon source coding.} The minimum average number
  of bits you need to encode a symbol drawn from a known
  distribution equals that distribution's entropy. Corollary: if
  you can predict, you can compress; if you can compress, you have
  (implicitly) predicted. Prediction and compression are the same
  job in different units.
  \item \textbf{Kolmogorov complexity.} The length of the shortest
  program that prints a string is, up to an additive constant that
  does not depend on the string, independent of which universal
  programming language you choose. So ``shortest description'' is
  a property of the string, not of the describer. This licenses
  talk of ``the'' complexity of a dataset.
  \item \textbf{Solomonoff induction.} If you must bet on what
  comes next with no prior knowledge, the Bayes-optimal gambler
  assigns probability proportional to $2^{-K}$, where $K$ is the
  Kolmogorov complexity of the hypothesis. Shorter explanations
  are more likely; this is Occam's razor made into a theorem.
  \item \textbf{AIXI.} An agent that acts to maximise expected
  reward under Solomonoff's prior is, in a precisely defined
  sense, the optimal reinforcement learner. It is uncomputable ---
  but every real learner is a computable approximation to it, and
  the distance from it is the distance from ideal intelligence.
\end{enumerate}

\begin{intuition}
\textbf{The triangle to keep on a napkin.} Three vertices ---
\emph{Predict}, \emph{Compress}, \emph{Act intelligently}. Three
edges. Shannon is the edge Predict\,$\leftrightarrow$\,Compress;
Solomonoff is the edge Compress\,$\leftrightarrow$\,Act (through
Kolmogorov, which makes ``compress'' language-independent); and the
full triangle closes at AIXI. \Cref{fig:triangle} draws it. If you
remember only one picture from these notes, remember that one.
\end{intuition}

\begin{figure}[h]
\centering
\begin{tikzpicture}[
  vertex/.style={draw, rounded corners, align=center, minimum width=2.6cm, minimum height=1cm, thick},
  edge/.style={-{Latex[length=2mm]}, thick}
]
  \node[vertex] (P) at (0,0) {Predict\\ \scriptsize next-token model};
  \node[vertex] (C) at (7,0) {Compress\\ \scriptsize shortest code};
  \node[vertex] (A) at (3.5,-4) {Act\\ \scriptsize reward-seeking agent};

  \draw[edge] (P) -- node[above, font=\small] {Shannon (\S2)} (C);
  \draw[edge] (C) -- node[right, font=\small, xshift=2pt] {Solomonoff (\S4)} (A);
  \draw[edge] (A) -- node[left, font=\small, xshift=-2pt] {AIXI (\S5)} (P);

  \node[font=\small\itshape] at (3.5,-1.8) {Kolmogorov (\S3)\\ \scriptsize makes the triangle\\ \scriptsize language-independent};
\end{tikzpicture}
\caption{The three-vertex picture. Each edge is a theorem that
translates one vertex into another. Kolmogorov complexity
(\S3) is the hub that makes the translation coordinate-free.}
\label{fig:triangle}
\end{figure}

\subsection*{Why this matters for decisions, not just for theory}

A decision-maker evaluating a foundation-model vendor faces a
simple but hard question: \emph{is the loss number on the vendor's
training dashboard a proxy for anything you care about?} If the
equivalence we are about to prove holds cleanly, the answer is yes:
every 0.1-bit improvement in held-out cross-entropy corresponds to
a measurable improvement in downstream capability, and the
conversion rate can be estimated. If the equivalence leaks --- and
we will show exactly where it does --- the answer is ``sometimes,
in these regimes, with these caveats.'' Either answer is
actionable; neither is available without the chain of theorems.

\begin{activity}[title={Activity A4 --- Falsification worksheet}]
Fill the table below before reading further. We will return to it
in \cref{sec:caveats} and see which of your rows survives.
\smallskip

\begin{tabular}{p{4.2cm}p{4.2cm}p{4.2cm}}
\toprule
\textbf{Equivalence claim} & \textbf{What would refute it} & \textbf{What is merely consistent} \\
\midrule
Lower bpc $\Rightarrow$ higher MMLU (same size). & & \\
Better zip ratio on a domain $\Rightarrow$ better agent on that domain. & & \\
A model that memorises (does not generalise) cannot compress held-out data well. & & \\
\bottomrule
\end{tabular}
\end{activity}

\paragraph{A caveat, promoted out of the footnotes.}
A \emph{perfect} next-token predictor would be an ideal learner.
Real next-token predictors are not perfect. The gap between them
and the ideal is exactly where this document's caveats live
(\cref{sec:caveats}); we state the thesis before we qualify it,
but the qualification is not optional. In particular: the
equivalence is between the \emph{training objective} and general
capability, not between any specific checkpoint and general
capability. A model can be a poor compressor of your domain while
the objective it was trained on remains the right objective. Keep
the distinction; a lot of vendor arguments blur it.

\medskip
\noindent\textbf{Formal statement of the thesis.}
The central technical claim of this note is that three operations
often treated as distinct --- predicting the next symbol in a
sequence, compressing that sequence into a shorter description, and
behaving intelligently in its presence --- are, under a precise set
of definitions, the \emph{same operation} viewed from three angles.
Shannon's source coding theorem \citep{shannon1948mathematical}
identifies prediction with compression; Kolmogorov complexity
\citep{kolmogorov1965three} makes the notion of ``shortest
description'' independent of the choice of description language;
Solomonoff's universal prior \citep{solomonoff1964formal} promotes
the Kolmogorov minimiser to a Bayesian predictor that converges to
any computable source; and Hutter's AIXI
\citep{hutter2005universal} upgrades that predictor into a
reward-seeking agent whose optimality can be stated as a theorem.
The chain is not an analogy. Every step in this note is a
definition, a lemma, or a theorem. The pedagogical payoff is that
one can then ask sharp questions about large language models ---
whether they count as compressors, whether better compression
entails more general capability, and where the equivalence breaks
--- without smuggling in metaphor.
